\documentclass{TDP003mall}

\usepackage{float}
\usepackage{minted}

\newcommand{\version}{Version 1.0}
\author{Otto Roming, \url{ottro821@student.liu.se}}
\title{Veckoplanering}
\date{2026-10-19}
\rhead{Otto Roming}

\begin{document}
\projectpage
\section{Versionshistorik}
\begin{table}[!h]
\begin{tabularx}{\linewidth}{|l|X|l|}
\hline
Ver. & Revisionsbeskrivning & Datum \\\hline
1.0 & Första version av veckoplanering & 2026-10-19 \\\hline
\end{tabularx}
\end{table}


\section{Bakgrund}
Just nu ligger jag bakom milstolpe-4 i kursen. Därför är det som viktigast att jag fokuserar på att bli klar med den innan halvtidsscreencasten så jag klarar av de obligatoriska delarna i kursen. Samtidigt satsar jag på betyg 5 i kursen vilket betyder att jag måste lyckas uppnå ett högt antal krav

\section{Vecka 43}
Målet denna vecka är att bli klar med sen inlämning av milstolpe-4.

\subsection{Tekniska krav}
\begin{itemize}
    \item Enkel inloggning av användare mot Firebase (1 p)
    \item Kontohantering, t.ex. lagra poäng, vänlista, eller dylikt som ska sparas över tid för olika användare. (1 p) 
\end{itemize}

\subsection{Entreprenöriella krav}
\begin{itemize}
    \item Firestore. Använder Cloud Firestore eller Realtime Database som databas. Båda lösningarna är Googles lagringstjänst för mobile- och webb-appar som erbjuder bl.a. databas och autentisering. Dessutom finns en mer generell databas-lösning som heter Firebase Cloud Storage. Cloud Storage accepterar både ljud, bild, och text så i princip alla appar kan utnyttja Cloud Storage på ett eller annat sätt. (2 p)
\end{itemize}

\subsection{Samanlagda poäng}
2 tekniska, 2 entreprenöriella


\section{Vecka 45}
Nu är det dags att göra färdigt halvtidsscreencast

\subsection{Tekniska krav}
\begin{itemize}
    \item Använder någon sensor, t.ex. kamera eller GPS (1p per sensor)
    \textit{2 poäng i mitt fall då jag använder mig av både GPS för att spara positionen av inlägg samt en kopplad bild}
    \item Använder forms och validerar användarinput, så att input är av rätt typ, och är korrekt. T.ex. att lösenord är tillräcklig starka. (2 p)
    \item Använder routes (go-routes biblioteket) för navigering i stället för Navigator.push. (1 p)    
\end{itemize}

\subsection{Entreprenöriella krav}
\begin{itemize}
    \item Sökfunktion i appen, dvs. att användaren kan skriva in ett sökord och få upp relevant information. Oftast krävs att sökningen genomförs på backend-sidan, för att säkerställa att all relevant information kan visas, inte bara det som råkade finnas på skrämen. (2 p)
\end{itemize}

\subsection{Samanlagda poäng}
7 tekniska, 4 entreprenöriella

\section{Vecka 46}
Implementera användarinställningar

\subsection{Tekniska krav}

\subsection{Entreprenöriella krav}
\begin{itemize}
    \item Någon typ av UX inställning, t.ex. möjlighet att ändra till dark mode (1 p)
    \item Någon form av accessibility stöd, t.ex. öka font-storlek. (1 p)
\end{itemize}

\subsection{Samanlagda poäng}
7 tekniska, 6 entreprenöriella

\section{Vecka 47}
Implementera en hjälpsam ai chattfunktion i appen

\subsection{Tekniska krav}

\subsection{Entreprenöriella krav}
\begin{itemize}
    \item Använder någon form av AI via publikt interface (API) (2 p)
\end{itemize}

\subsection{Samanlagda poäng}
7 tekniska, 8 entreprenöriella

\section{Vecka 48 och 49}
Finslipa koden till en komplett applikation

\subsection{Tekniska krav}
\begin{itemize}
    \item{Hantering av stora och små skärmar, eller skärmrotation, genom att använda olika layouter. Användaren ska kunna byta från mobil till tablet utan att tappa i upplevelse. Bl.a. är det viktigt att information användaren hunnit skriva in bevaras, så att skärmen förblir detsamma. Det är också viktigt att layout:en ser bra ut i båda lägen. Användarupplevelsen kan enkelt testas genom att rotera mobilen eller emulatorn före uppstart. (2 p)}
\end{itemize}

\subsection{Entreprenöriella krav}

\subsection{Samanlagda poäng}
9 tekniska, 8 entreprenöriella

\subsection{Samanlagda poäng}
7 tekniska, 8 entreprenöriella

\section{Vecka 50 och 51}
\begin{itemize}
    \item slutdemo
    \item inlämning av slut-screencast på ca. 5 minuter
    \item inlämning av kod
\end{itemize}

\subsection{Tekniska krav}

\subsection{Entreprenöriella krav}

\subsection{Samanlagda poäng}
9 tekniska, 8 entreprenöriella


% \bibliographystyle{IEEEtran}
% \bibliography{references}

\end{document}